\section{Overview}

\subsection{The research question}

Do structurally different \emph{regime processes}, used as the gate of a mixture of experts, change how
well a model predicts which US stocks will outperform which over the next week? The thesis compares
gates, not network architectures: the experts, the base network, the data and the training protocol
are held fixed, and only the process that produces the regime weights varies across arms.

\dec{2026-10-08, Q2 and Q4} The thesis is \textbf{comparison-led}: ``which regime mechanism makes the best
gate''. No constructive element (such as a new differentiable gate) is added.

\begin{whybox}
\begin{itemize}
\item \textbf{The mechanism itself is not new, the comparison is.} A novelty audit found that gated experts for
  regime discovery go back to Weigend, Mangeas \& Srivastava (1995), that an HMM has been used as a
  mixture-of-experts gate (2022), and that filtering-based gating on financial data is published (MoE-F,
  ICLR 2025). What nobody has done is a \emph{controlled} comparison of structurally different regime
  processes (Markov switching, jump models, distributional clustering, covariate-driven switching) as the
  gate of the same model on the same data. That is a genuine gap, and it is the thesis.
\item \textbf{A comparison is informative whichever gate wins.} A design-led thesis (``here is a new gate and it
  beats the others'') succeeds only if the new gate wins. A comparison produces a result in every case,
  including ``the regime process does not matter'', which is itself useful to practitioners who must pick
  one.
\item \textbf{It fits the time and compute available.} A new differentiable gate would need its own
  derivation, implementation and validation on top of the comparison.
\end{itemize}
\giveup{a ``new method'' headline. Positioning against the closest precedent (Ye \& Borde) is still open
(Q5).}
\end{whybox}

\subsection{The model in one line}

For stock $i$ on date $t$, with characteristic vector $x_{i,t}$,
\begin{equation}
\hat y_{i,t} \;=\; \underbrace{f_0(x_{i,t})}_{\text{frozen base}} \;+\; \sum_{k=1}^{K}
\underbrace{\pi_k(i,t)}_{\text{frozen gate}}\;\underbrace{r_k(x_{i,t})}_{\text{expert correction}} .
\label{eq:model}
\end{equation}

\begin{itemize}
\item $f_0$ is the \textbf{base}: a multilayer perceptron (MLP, a feed-forward neural network) trained
  first, then frozen. It captures the regime-independent relation between characteristics and returns.
\item $r_k$ are $K$ \textbf{experts}: small MLPs whose outputs start at zero and learn only
  \emph{corrections} to the base (the residual design).
\item $\pi_k(i,t)$ are the \textbf{gate weights}: probabilities of the regimes, produced by a regime model
  fitted separately on market (and, now, industry) series and then frozen. Before 8 October the
  weight was the same for every stock on a date; it is now partly stock-specific through the stock's
  industry (Q1, Section~\ref{sec:industrygate}).
\end{itemize}

\subsection{Design principles and why}

\paragraph{The gate is fitted separately and frozen.} \dec{2026-09-21, Q19} Every gate is fitted on its own
data (the market series) before the experts are trained, and its weights are then fixed inputs.

\begin{whybox}
\begin{itemize}
\item \textbf{We want to see how the market regime itself affects the cross-section of returns}, even though
  training the gate together with the experts would give a better-optimised model. The question is
  ``when the market is in state $k$, how does the relation between characteristics and returns change?''.
  That question needs a regime that is defined by the market, not by what helps the loss.
\item \textbf{Joint training would turn the gate into a learned router.} If the gate received gradients from
  the experts' loss, it would move its boundaries to wherever they reduce prediction error, whether or not
  those boundaries correspond to market states. Any improvement could then come from the extra capacity
  and the freedom to place boundaries, not from the regime structure, and we could no longer attribute it.
\item \textbf{It is what makes the gates comparable.} Under joint training every gate type would be reshaped
  by the same gradient and the differences between them would wash out. Frozen, every arm receives a
  regime assignment it cannot influence, and the only thing that differs across arms is the process that
  produced it. That is a controlled experiment.
\item \textbf{Two of the gates cannot be trained by backpropagation at all} (the jump model and Wasserstein
  clustering are fitted by discrete optimisation). Joint training would favour the differentiable gates
  for a reason unrelated to the question.
\item \textbf{The regimes stay interpretable.} A frozen gate's regimes are calm and stressed market states that
  can be checked against history (2008, 2020). A jointly trained gate's ``regimes'' have no fixed meaning.
\end{itemize}
\giveup{the best achievable fitted loss. Jointly trained routing does better on the training criterion
(the DeepSeek mixture-of-experts work makes this point). We keep jointly trained gates as baseline arms,
so the gap between a co-adapted partition and the frozen ones is reported, not hidden.}
\end{whybox}

\paragraph{The base is pre-trained and frozen.} \dec{2026-09-21, Q7}

\begin{whybox}
\begin{itemize}
\item \textbf{Same logic as the gate.} We measure what regime-conditional corrections add to a fixed
  baseline. If the base kept learning during expert training it could reorganise itself around the
  mixture, and part of any gain would be the base improving, not the regimes helping.
\item \textbf{It defines the headline number.} The thesis reports the \emph{improvement over the base}. That
  number only means something if the base is the same object in every arm.
\item \textbf{It is shared across all gates.} The base never sees regime information, so one base per fold and
  seed serves every gate arm: a large compute saving, and a guarantee that every gate is compared on top
  of an identical foundation.
\end{itemize}
\giveup{the configuration likely to fit best (base and experts trained together, as in DeepSeekMoE's
shared expert). It is kept as construction arm B, so its effect is still measured.}
\end{whybox}

\paragraph{Residual mixture.} The experts correct a base instead of each producing a full forecast.

\begin{whybox}
\begin{itemize}
\item \textbf{It isolates the gate's contribution in one term}, $\sum_k\pi_k r_k$, which can be measured,
  plotted through time, and tested for concentration at regime transitions. In a standard mixture the
  gate's effect is tangled up with each expert's full forecast.
\item \textbf{The signal is weak.} Predictable variation in weekly returns is a fraction of a percent of
  their variance. Learning a small correction to an already good forecaster is a much easier problem, with
  less estimation noise, than learning $K$ full forecasters from scratch.
\item \textbf{Evidence from the closest precedent} (Ye \& Borde): a frozen base with zero-initialised experts
  was more accurate and far more stable than a standard mixture (0 of 30 seeds collapsed against 24 of 30).
  Their stability measure is volatility-specific, so the transfer to returns is checked, not assumed
  (Q6).
\end{itemize}
\giveup{experts cannot replace the base forecast entirely in a regime where the base is badly wrong; they
can only move it.}
\end{whybox}

\paragraph{Report the full mechanism-by-depth grid.} \dec{2026-09-19, Q11}

\begin{whybox}
Which expert depth is best may depend on the gate: a sharp gate gives each expert fewer effective
observations and favours smaller experts, a diffuse gate the opposite. Fixing one depth for all gates could
handicap some; tuning depth per gate would break the control. Reporting the whole grid contains the
fixed-depth comparison as one column and shows the interaction instead of hiding it. If the ranking of
gates is the same at every depth, the conclusion is robust; if it changes, that is a finding.
\giveup{compute (one full comparison per depth) and a larger multiple-testing burden (Q15).}
\end{whybox}

\paragraph{No width sweep and no sweep over $K$.} \dec{2026-09-18, advisor meeting}

\begin{whybox}
The advisor's point was scope: everything as first written was impossible, and some choices must be
made by argument. Widths follow a pyramid rule from one first width, since the first layer carries most of
the parameters anyway. $K$ is fixed by argument and held across all gates, because a sweep over $K$
multiplies cost and multiple testing, and because the data, not a search, limit how many regimes can be
estimated (Q18).
\giveup{the chance that a different width or $K$ would favour some gate; this is acknowledged, not tested.}
\end{whybox}

\paragraph{Research integrity.} No out-of-sample number on the real panel is computed before the design
is pre-registered; licensed data never leave \texttt{Data/}.

\begin{whybox}
The effects studied are small, and the design has many arms. If test results could influence choices,
selection noise could exceed the true differences between gates (Q15). Fixing everything on pre-2010 data
and writing it down before scoring is the only reliable protection.
\end{whybox}

\subsection{The pipeline at a glance}

\begin{center}
\resizebox{\textwidth}{!}{%
\begin{tikzpicture}[node distance=5mm and 6mm, box/.style={draw,rounded corners,align=center,
  minimum width=27mm,minimum height=9mm,font=\small}, >=Latex]
\node[box,fill=gray!10] (crsp) {CRSP daily\\(prices, quotes,\\membership)};
\node[box,fill=gray!10,below=of crsp] (comp) {Compustat\\(fundamentals,\\GICS, link)};
\node[box,right=12mm of crsp,yshift=-7mm] (panel) {Point-in-time\\panel: 57 ranked\\inputs, target};
\node[box,above=of panel] (mkt) {Market index\\daily log return};
\node[box,right=of panel] (base) {Base $f_0$\\(MLP, frozen)};
\node[box,right=of mkt] (gate) {Gate: market \&\\industry regimes\\(frozen)};
\node[box,right=of base,yshift=7mm] (exp) {Experts $r_k$\\(trained)};
\node[box,right=of exp] (eval) {Walk-forward\\evaluation};
\draw[->] (crsp) -- (panel); \draw[->] (comp) -- (panel);
\draw[->] (crsp.north) |- (mkt);
\draw[->] (panel) -- (base); \draw[->] (mkt) -- (gate);
\draw[->] (base) -- (exp); \draw[->] (gate) -| (exp);
\draw[->] (panel.south) -- ++(0,-6mm) -| (exp.south);
\draw[->] (exp) -- (eval);
\end{tikzpicture}}
\end{center}

Order of work set by the advisor (18 September): data, then compute budget, then parameters; start
experimenting with the base and a generic loss before refining the loss. The reason: the architecture
choices can only be settled once the compute budget is calculated rather than guessed, and the budget
depends on the data.
